
% ════════════════════════════════════════════
% 模型假设（03.模型假设.tex）—— 写作提示词见本文件末尾【本章节写作提示词】。
% ════════════════════════════════════════════
\section{模型假设}

为简化问题，本文做出以下假设：

\begin{itemize}[itemindent=2em]
\item \textbf{假设1：} 观测车在问题一至问题三中视为具有等效半径的圆盘且匀速行驶，忽略加速、减速、转向与避碰时间，从而可用几何距离与恒定速度精确刻画车辆时间轴。
\item \textbf{假设2：} 方位角测量误差、车辆定位误差、云台零位偏差与控制误差均为相互独立的零均值正态随机变量，保证线性误差传播与方差叠加公式成立。
\item \textbf{假设3：} 视线与建筑物棱柱的棱、面或顶点相切或接触即视为遮挡，采用严格的包含边界判交准则，保证可观测区域判定偏于安全。
\item \textbf{假设4：} 所有无人机目标在部署阶段视为静止，且目标高度由告警信息直接给出而不参与交会解算，从而将三维定位问题分解为水平交会与俯仰约束两部分。
\end{itemize}
